\documentclass[10pt,a4paper]{amsart}
\usepackage[margin=19mm]{geometry}
\usepackage{amsmath,amssymb,mathtools}
\usepackage[scr=rsfs,cal=euler]{mathalfa}
\usepackage{xfrac}
\usepackage[unicode,hidelinks,bookmarks=false]{hyperref}
\newcommand{\ie}{i.e.\ }
\newcommand{\cf}{cf.\ }
\newcommand{\resp}{respectively\ }
\newcommand{\Subsection}{\subsection}
\providecommand{\nobreakdash}{\nobreak\hbox{-}}
\setlength{\emergencystretch}{2em}
\multlinegap0pt
\usepackage[shortlabels]{enumitem}
\usepackage{xcolor}
\usepackage{amssymb}
\usepackage{mathtools}
\newtheorem{lem}{Lemma}
\newcommand{\R}{\mathbb{R}}
\title{Two-boost problem\\ for the Newtonian potential at the infinity}
\author{Jagna Wi\'sniewska}
\date{Source: arXiv:2512.13373v2 (CC BY 4.0)}
\begin{document}
\maketitle

\noindent\emph{Source and licence.} Two-boost problem\\ for the Newtonian potential at the infinity. Version arXiv:2512.13373v2 (CC BY 4.0), distributed by arXiv under the Creative Commons Attribution 4.0 International licence, http://creativecommons.org/licenses/by/4.0/, as recorded in the arXiv licence metadata for this exact version and shown on the abstract page https://arxiv.org/abs/2512.13373v2. Full licence notice: \texttt{licenses/arxiv-2512.13373-CC-BY-4.0.txt}.\par\smallskip

\noindent\textit{Editorial scope.} Counted unit: the article's lem lem:noMax (source0.tex lines 416--422). The article's complete original proof of it is delivered with it (source0.tex lines 423--450, one \texttt{proof} environment of 28 lines). No mathematics is written, repaired or re-derived: the statement and the proof are the article's own lines, in the article's own notation, and no retained line is substituted or re-flowed.  Retained uncounted as the source-local context the proof needs: lines 102--104; lines 106--108; lines 117--119; lines 184--189; lines 294--298; lines 370--377.  Nothing was omitted from any retained span, so the package carries no omission record.  No retained line contains a citation, so the wrapper carries no bibliography and no bibliography entry.  No retained line contains a figure environment, an image-inclusion command or drawing code, so no image is delivered and the package declares no asset dependency.  Editorial formatting only, in addition: an AMS \texttt{amsart} preamble that restates the article's own declarations and macros used by the retained lines, namely the environments \texttt{lem} on separate counters, together with the standard packages those lines need.  The retained environments print in this document as Lemma 1 in order of appearance, whereas the article prints the same items with the numbering of its own full document.  The complete native source of the article as distributed in the arXiv e-print for this exact version is bundled unchanged as \texttt{sources/191022/source0.tex}.
\begin{equation}\label{DefH}
H(q,p):= H_0(q,p) - V(q),
\end{equation}

\begin{equation}\label{DefV}
V(r, \theta) \leq \frac{a}{r}\qquad\textrm{and}\qquad \partial_r V(r,\theta) +\frac{2}{r}V(r,\theta)\geq 0.
\end{equation}

\begin{equation}\label{CondC}
c> \max\left\lbrace \sqrt[3]{32a^2}, \sqrt{4a\left(3 R_1+2\sqrt[3]{2a}\right)}\right\rbrace,
\end{equation}

\begin{lem}\label{lem:energy}
Let $H$ be the Hamiltonian as in \eqref{DefH} and choose an energy value $c> \sqrt{2aR_1^2} $. Then for all $(r,\theta, p_r, p_\theta)\in H^{-1}(c)$, such that $r\in [R_1, \frac{c^2}{2a})$, we have
$$
c-p_\theta>\frac{c^2-2ar}{2(c+r^2)}.
$$
\end{lem}

\begin{align}
\chi_0(r) & :=\chi \left(\frac{r-R_1}{R_2-R_1}\right), \label{DefChi0}\\
\chi_1(q,p) & := \chi(H_0(q,p)-\sup V-c),\label{DefChi1}\\
H_1 &:= H_0 - \chi_0\chi_1 V.\label{DefH1}
\end{align}

\begin{align}
\{H_1, \{H_1, r\}\}(x) & =\frac{p_\theta^2}{r^3}+\chi_0\partial_rV+\chi_0' V\nonumber \\
& = \frac{2}{r}\left( \frac{p_\theta^2}{2r^2}-\chi_0 V\right)+\chi_0\left( \frac{2}{r}V+ \partial_rV\right)+\chi_0' V\nonumber\\
& = \frac{2}{r}(H_1-p_\theta)+\chi_0\left( \frac{2}{r}V+ \partial_rV\right)+\chi_0' V\nonumber\\
& = \frac{2}{r}(H_1-p_\theta)+\chi_0\left( \frac{2}{r}V+ \partial_rV\right)-\frac{2}{R_2-R_1} V\nonumber\\
& \geq  \frac{2}{r}\left( H_1-p_\theta-\frac{a}{R_2-R_1}\right)\label{H1H1r}\\
& \geq \frac{2}{r}\left( e-\frac{a}{R_2-R_1}\right)=\frac{e}{r}>0,\nonumber
\end{align}

\begin{lem}\label{lem:noMax}
Consider a potential function as in \eqref{DefV}. Choose an energy value $c> 0$ as in \eqref{CondC}.
Then for a Hamiltonian $H_1$ as in \eqref{DefH1} with $R_2:= \frac{1}{8a}\left( c^2+2aR_1\right)$ we have
$$
H_1^{-1}(c)\cap\{\{H_1,r\}=0\}\cap \{\{H_1,\{H_1,r\}\}\leq 0\}\cap T^*(B(R_2)\setminus B(R_1))=\emptyset.
$$
\end{lem}

\begin{proof}
By assumption \eqref{CondC} we have
\begin{align}
R_2 & = \frac{1}{8a}\left( c^2+2aR_1\right) \geq \frac{1}{8a}\left( 12aR_1+2aR_1\right)=\frac{7}{4} R_1 \nonumber\\
R_2 & = \frac{1}{8a}\left( c^2+2aR_1\right)\leq\frac{1}{8a}\left( c^2+\frac{1}{6}c^2\right)=\frac{7c^2}{48 a}<\frac{c^2}{2a}.\label{R2small}
\end{align}

Moreover, observe that for every $x=(r,\theta, p_r, p_\theta) \in T^*(\R^2\setminus B(R_1))\cap H_1^{-1}(c)$ we have
$$
c-p_\theta \geq \frac{p_\theta^2}{2 r^2}-\chi_0(r)V(r,\theta)\geq \frac{p_\theta^2}{2 r^2}-\frac{a}{r}.
$$
Consequently, by a similar argument as the one presented in the proof of Lemma \ref{lem:energy}, we can infer that for all $x\in H^{-1}(c)\cap T^*(B(R_2)\setminus B(R_1))$ we have
$$
c-p_\theta \geq \frac{c^2-2aR_2}{2(c+R_2^2)}.
$$
Now we can use relation \eqref{H1H1r}, the inequalities above, and assumption\linebreak $R_2=\frac{1}{8a}\left( c^2+2aR_1\right)$ to estimate
\begin{align*}
\{H_1, \{H_1, r\}\}(x) & \geq  \frac{2}{r}\left( c-p_\theta-\frac{a}{R_2-R_1}\right)\\
& \geq \frac{2}{r}\left(\frac{c^2-2aR_2}{2(c+R_2^2)}-\frac{a}{R_2-R_1}\right)\\
& = \frac{-4a R_2^2+R_2(c^2+2aR_1)-c(2a+cR_1)}{r(c+R_2^2)(R_2-R_1)}\\
& = \frac{\frac{1}{16a}\left( c^2+2aR_1\right)^2-c(2a+cR_1)}{r(c+R_2^2)(R_2-R_1)}\\
& = \frac{c^2\left(c^2-12aR_1\right)-32a^2c+4a^2R_1^2}{16ar(c+R_2^2)(R_2-R_1)}\\
& \geq \frac{8ac^2\sqrt[3]{2a}-32a^2c+4a^2R_1^2}{16ar(c+R_2^2)(R_2-R_1)}\\
& = \frac{2c\left(c\sqrt[3]{2a}-4a\right)+aR_1^2}{4r(c+R_2^2)(R_2-R_1)}\\
& \geq \frac{aR_1^2}{4r(c+R_2^2)(R_2-R_1)}>0.
\end{align*}
The estimates in the last three lines come from the conditions \eqref{CondC} on the energy value $c$.
\end{proof}

\end{document}
